\documentclass[11pt, latterpaper]{article}
\usepackage [utf8]{inputenc}
\usepackage [resetlabels, labeled]{multibib}
\usepackage [T2A]{fontenc}
\usepackage [top=2cm, bottom=2cm, left=1.5cm, right=1.5cm]{geometry}
\usepackage {graphicx}
\usepackage{amsmath}
\usepackage{pgfplots}

\graphicspath{{/}}

\linespread{1.5} %interval

\title{KaTeX document title}
\author{lyalya \thanks{nyanya}}
\date{21.09.2026}
\begin{document}
\maketitle
We have now added a \textbf{title}, \underline{author} and \textit{date} to our first \LaTeX document!

Ляляля
\section{Первый блок документа}

\section{Формулы триганометрии}
\subsection{Основное тригонометрическое тождество}
 $$sin^2 \alpha + cos^2 \alpha = 1 $$
\subsection{Формулы сложения}
$$sin(\alpha + \beta)=sin\alpha*cos\beta+cos\alpha*sin\beta$$
\subsection{Формулы двойного угла}
$$sin2\alpha=2sin\alpha cos\alpha$$
\subsection{Формулы понижения степени}
$$sin^2\alpha=\frac{1-cos2\alpha}{2}$$
\section{Графики функций}
\subsection{Синуса}
\begin{center}
\begin{tikzpicture}
\begin{axis}[
	title=Синус,
	xlabel={x},
	ylabel={y},
	domain=-6.28:6.28, % Область от -2pi до 2pi
	samples=100,        % Сглаживание графика
	minor tick num=2,
	axis lines=middle   % Пересечение осей в центре
]
\addplot[blue, thick] {sin(deg(x))}; % deg(x) преобразует радианы в градусы для корректного расчета
\end{axis}
\end{tikzpicture}
\end{center}

\subsection{Косинуса}
\begin{center}
\begin{tikzpicture}
\begin{axis}[
	title=Косинус,
	xlabel={x},
	ylabel={y},
	domain=-6.28:6.28,
	samples=100,
	minor tick num=2,
	axis lines=middle
]
\addplot[red, thick] {cos(deg(x))};
\end{axis}
\end{tikzpicture}
\end{center}

\section{Таблицы}

\subsection{Синусы}
Полноразмерная таблица значений функции $y = \sin \alpha$:
\begin{center}
\renewcommand{\arraystretch}{2} % Увеличиваем вертикальные отступы в таблице для дробей
\begin{tabular}{|c|c|c|c|c|c|c|c|c|c|}
\hline
$\alpha$ (град.) & $0^\circ$ & $30^\circ$ & $45^\circ$ & $60^\circ$ & $90^\circ$ & $120^\circ$ & $135^\circ$ & $150^\circ$ & $180^\circ$ \\ \hline
$\alpha$ (рад.)  & $0$ & $\dfrac{\pi}{6}$ & $\dfrac{\pi}{4}$ & $\dfrac{\pi}{3}$ & $\dfrac{\pi}{2}$ & $\dfrac{2\pi}{3}$ & $\dfrac{3\pi}{4}$ & $\dfrac{5\pi}{6}$ & $\pi$ \\ \hline
$\sin \alpha$   & $0$ & $\dfrac{1}{2}$ & $\dfrac{\sqrt{2}}{2}$ & $\dfrac{\sqrt{3}}{2}$ & $1$ & $\dfrac{\sqrt{3}}{2}$ & $\dfrac{\sqrt{2}}{2}$ & $\dfrac{1}{2}$ & $0$ \\ \hline
\end{tabular}
\end{center}

\subsection{Косинусы}
Полноразмерная таблица значений функции $y = \cos \alpha$:
\begin{center}
\renewcommand{\arraystretch}{2}
\begin{tabular}{|c|c|c|c|c|c|c|c|c|c|}
\hline
$\alpha$ (град.) & $0^\circ$ & $30^\circ$ & $45^\circ$ & $60^\circ$ & $90^\circ$ & $120^\circ$ & $135^\circ$ & $150^\circ$ & $180^\circ$ \\ \hline
$\alpha$ (рад.)  & $0$ & $\dfrac{\pi}{6}$ & $\dfrac{\pi}{4}$ & $\dfrac{\pi}{3}$ & $\dfrac{\pi}{2}$ & $\dfrac{2\pi}{3}$ & $\dfrac{3\pi}{4}$ & $\dfrac{5\pi}{6}$ & $\pi$ \\ \hline
$\cos \alpha$   & $1$ & $\dfrac{\sqrt{3}}{2}$ & $\dfrac{\sqrt{2}}{2}$ & $\dfrac{1}{2}$ & $0$ & $-\dfrac{1}{2}$ & $-\dfrac{\sqrt{2}}{2}$ & $-\dfrac{\sqrt{3}}{2}$ & $-1$ \\ \hline
\end{tabular}
\end{center}

\subsection{Арксинусы}
Таблица главных значений обратной тригонометрической функции $y = \arcsin x$:
\begin{center}
\renewcommand{\arraystretch}{2}
\begin{tabular}{|c|c|c|c|c|c|c|c|}
\hline
$x$ & $-1$ & $-\dfrac{\sqrt{3}}{2}$ & $-\dfrac{1}{2}$ & $0$ & $\dfrac{1}{2}$ & $\dfrac{\sqrt{3}}{2}$ & $1$ \\ \hline
$\arcsin x$ & $-\dfrac{\pi}{2}$ & $-\dfrac{\pi}{3}$ & $-\dfrac{\pi}{6}$ & $0$ & $\dfrac{\pi}{6}$ & $\dfrac{\pi}{3}$ & $\dfrac{\pi}{2}$ \\ \hline
\end{tabular}
\end{center}

\subsection{Арккосинусы}
Таблица главных значений обратной тригонометрической функции $y = \arccos x$:
\begin{center}
\renewcommand{\arraystretch}{2}
\begin{tabular}{|c|c|c|c|c|c|c|c|}
\hline
$x$ & $-1$ & $-\dfrac{\sqrt{3}}{2}$ & $-\dfrac{1}{2}$ & $0$ & $\dfrac{1}{2}$ & $\dfrac{\sqrt{3}}{2}$ & $1$ \\ \hline
$\arccos x$ & $\pi$ & $\dfrac{5\pi}{6}$ & $\dfrac{2\pi}{3}$ & $\dfrac{\pi}{2}$ & $\dfrac{\pi}{3}$ & $\dfrac{\pi}{6}$ & $0$ \\ \hline
\end{tabular}
\end{center}

\textit{\underline{\textbf{Some text}}}
\end{document}